\begin{proof}[Proof of \cref{obj:thm:five-output-region}]
\begin{enumerate}
\item \textit{Certificate coordinates and profiling equations.}
Let \(\mathcal K_5\) be the active-pattern set and let \(\mathcal U\) be the interior parameter domain:
\[
\mathcal K_5=\{S_3,S_4,S_6,S_8,S_9\},
\qquad
\mathcal U=(0,1)^3\times(0,\infty).
\]
For \(\xi=(p,\mu_0,\mu_1,\varepsilon)\in\mathcal U\), set
\[
\theta(\xi)=(\mu_0,\mu_1)^\top,\qquad
q=1-p,\qquad
r_\varepsilon=e^\varepsilon,\qquad d_\varepsilon=r_\varepsilon-1.
\]
Use the pattern masses, gradients, profiling directions, and projected scores from
\cref{obj:lem:staircase-refinement}. Write their dependence on \(\xi\) explicitly as
\[
h_S(\xi)=h_S(\theta(\xi)),\qquad
g_S(\xi)=\nabla_\theta h_S(\theta(\xi)),\qquad
f_S(\xi,t)=\frac{(g_S(\xi)^\top v(t))^2}{h_S(\xi)},
\qquad
\rho_S(\xi,t)=\frac{g_S(\xi)^\top v(t)}{h_S(\xi)}.
\]
Here \(p,\varepsilon\) are fixed in the gradient. In particular,
\[
v(t)=\begin{pmatrix}t\\t+1\end{pmatrix},
\qquad
e_{\mathrm{prof}}=\begin{pmatrix}1\\1\end{pmatrix},
\qquad
F_\theta(\alpha,t)=\sum_{S\in\mathcal S}\alpha_S f_S(\xi,t).
\]

Index the four records by \(x_0,x_1,x_2,x_3\). Their probabilities, in this order, are
\[
(\pi_{\theta,0},\pi_{\theta,1},\pi_{\theta,2},\pi_{\theta,3})
=
\bigl(q(1-\mu_0),q\mu_0,p(1-\mu_1),p\mu_1\bigr).
\]
The pattern enumeration is specified by
\[
x_j\in S_m
\quad\Longleftrightarrow\quad
\left\lfloor\frac{m}{2^j}\right\rfloor\equiv1\pmod 2,
\qquad
0\le j\le3,\quad 1\le m\le14.
\]
Define the ray vectors and ray matrix by
\[
(b_S(\varepsilon))_j=
\begin{cases}
r_\varepsilon,&x_j\in S,\\
1,&x_j\notin S,
\end{cases}
\qquad
B_\varepsilon=(b_{S_1}(\varepsilon),\ldots,b_{S_{14}}(\varepsilon)).
\]
All record probabilities and ray coordinates are positive, so
\[
h_S(\xi)=\sum_{j=0}^3\pi_{\theta,j}(b_S(\varepsilon))_j>0.
\]

The stationarity, dual, and score conditions described in
\cref{obj:def:support-regions} enter the following full certificate definition.
A point \(\xi\in\mathbb R^4\) belongs to \(\mathcal R_5\) if and only if
\(\xi\in\mathcal U\) and there exist
\(\alpha\in\mathbb R^{14}\), \(t\in\mathbb R\), and \(\eta\in\mathbb R^4\) satisfying
\[
\begin{gathered}
\alpha\ge0,\qquad B_\varepsilon\alpha=\mathbf1,\qquad
\alpha_S>0\ (S\in\mathcal K_5),\qquad
\alpha_S=0\ (S\notin\mathcal K_5),\\
\partial_tF_\theta(\alpha,t)=0,\\
b_S(\varepsilon)^\top\eta=f_S(\xi,t)\quad(S\in\mathcal K_5),\\
f_S(\xi,t)<b_S(\varepsilon)^\top\eta\quad(S\notin\mathcal K_5),\\
\rho_S(\xi,t)\ne\rho_T(\xi,t)
\quad(S,T\in\mathcal K_5,\ S\ne T).
\end{gathered}
\]
The derivative appearing here is
\[
\partial_tF_\theta(\alpha,t)
=
2\sum_S\alpha_S
\frac{(g_S(\xi)^\top v(t))(g_S(\xi)^\top e_{\mathrm{prof}})}
     {h_S(\xi)}.
\]

The five active columns of the ray matrix are
\[
B_\varepsilon\big|_{\mathcal K_5}
=
\begin{pmatrix}
r_\varepsilon&1&1&1&r_\varepsilon\\
r_\varepsilon&1&r_\varepsilon&1&1\\
1&r_\varepsilon&r_\varepsilon&1&1\\
1&1&1&r_\varepsilon&r_\varepsilon
\end{pmatrix},
\]
with columns ordered as \(S_3,S_4,S_6,S_8,S_9\).
Subtracting consecutive normalization equations, and using
\(r_\varepsilon-1>0\), gives
\[
\alpha_{S_3}=\alpha_{S_4}=\alpha_{S_8},\qquad
\alpha_{S_6}=\alpha_{S_9},\qquad
(r_\varepsilon+2)\alpha_{S_3}
+(r_\varepsilon+1)\alpha_{S_6}=1.
\]
Thus every feasible vector supported within \(\mathcal K_5\) has the form
\[
\alpha^u_S(\varepsilon)=
\begin{cases}
u/(r_\varepsilon+1),&S\in\{S_6,S_9\},\\
(1-u)/(r_\varepsilon+2),&S\in\{S_3,S_4,S_8\},\\
0,&S\notin\mathcal K_5,
\end{cases}
\qquad 0\le u\le1.
\]
Conversely, the displayed matrix verifies feasibility of this vector.
Its five active weights are positive precisely when \(0<u<1\).

Define the two normalized constituent objectives and their coefficients by
\[
\begin{aligned}
\Phi_{\mathrm{bin}}(\xi,t)
&=\frac{f_{S_6}(\xi,t)+f_{S_9}(\xi,t)}{r_\varepsilon+1},\\
\Phi_{\mathrm{ter}}(\xi,t)
&=\frac{f_{S_3}(\xi,t)+f_{S_4}(\xi,t)+f_{S_8}(\xi,t)}
        {r_\varepsilon+2},\\
a_{\mathrm{bin}}(\xi)
&=\frac{d_\varepsilon^2}{r_\varepsilon+1}
  \left(\frac1{h_{S_6}(\xi)}+\frac1{h_{S_9}(\xi)}\right),\\
a_{\mathrm{ter}}(\xi)
&=\frac{d_\varepsilon^2p^2}{r_\varepsilon+2}
  \left(\frac1{h_{S_4}(\xi)}+\frac1{h_{S_8}(\xi)}\right).
\end{aligned}
\]
Both coefficients are positive throughout \(\mathcal U\).
The active gradients are
\[
g_{S_3}=0,\quad
g_{S_4}=\begin{pmatrix}0\\-d_\varepsilon p\end{pmatrix},\quad
g_{S_6}=\begin{pmatrix}d_\varepsilon q\\-d_\varepsilon p\end{pmatrix},\quad
g_{S_8}=\begin{pmatrix}0\\d_\varepsilon p\end{pmatrix},\quad
g_{S_9}=-g_{S_6}.
\]
Consequently,
\[
\begin{aligned}
\Phi_{\mathrm{bin}}(\xi,t)
&=a_{\mathrm{bin}}(\xi)\{p-(1-2p)t\}^2,\\
\Phi_{\mathrm{ter}}(\xi,t)
&=a_{\mathrm{ter}}(\xi)(t+1)^2,\\
F_\theta(\alpha^u,t)
&=u\Phi_{\mathrm{bin}}(\xi,t)+(1-u)\Phi_{\mathrm{ter}}(\xi,t).
\end{aligned}
\]
The binary rays sum to \((r_\varepsilon+1)\mathbf1\), and the ternary rays
sum to \((r_\varepsilon+2)\mathbf1\). Averaging the active dual equalities
over these two families therefore yields
\[
\Phi_{\mathrm{bin}}(\xi,t)
=
\sum_{j=0}^3\eta_j
=
\Phi_{\mathrm{ter}}(\xi,t).
\]
Finally, the certificate stationarity equation becomes
\[
u\,\partial_t\Phi_{\mathrm{bin}}(\xi,t)
+(1-u)\,\partial_t\Phi_{\mathrm{ter}}(\xi,t)=0.
\]
% lean: CausalSmith.Stat.LdpAteEfficiencySurface.r5_stationarity_eq_mixture

\item \textit{Continuous continuation and openness.}
Fix
\[
\xi_\circ=(p_\circ,\mu_{0,\circ},\mu_{1,\circ},\varepsilon_\circ)
\in\mathcal R_5
\]
and a certificate \((\alpha_\circ,t_\circ,\eta_\circ)\).
Define, for \(\xi\in\mathcal U\),
\[
s_{\mathrm{loc}}(\xi)
=
\sqrt{\frac{a_{\mathrm{ter}}(\xi)}{a_{\mathrm{bin}}(\xi)}}.
\]
The constituent equality implies
\[
\{p_\circ-(1-2p_\circ)t_\circ\}^2
=
\{s_{\mathrm{loc}}(\xi_\circ)(t_\circ+1)\}^2.
\]
Splitting equality of squares gives a sign
\(\sigma\in\{-1,1\}\) such that
\[
p_\circ-(1-2p_\circ)t_\circ
=
\sigma s_{\mathrm{loc}}(\xi_\circ)(t_\circ+1).
\]
Its rearrangement is
\[
\{1-2p_\circ+\sigma s_{\mathrm{loc}}(\xi_\circ)\}t_\circ
=
p_\circ-\sigma s_{\mathrm{loc}}(\xi_\circ).
\]
If the coefficient on the left were zero, the right side would also be
zero, giving \(p_\circ=1\). Thus, on a neighborhood where its denominator
is nonzero, define
\[
t_\sigma(\xi)
=
\frac{p-\sigma s_{\mathrm{loc}}(\xi)}
     {1-2p+\sigma s_{\mathrm{loc}}(\xi)}.
\]
This function is continuous at \(\xi_\circ\), equals \(t_\circ\) there,
and satisfies
\[
\Phi_{\mathrm{bin}}(\xi,t_\sigma(\xi))
=
\Phi_{\mathrm{ter}}(\xi,t_\sigma(\xi)).
\]

At the original certificate,
\[
\partial_t\Phi_{\mathrm{ter}}(\xi_\circ,t_\circ)
=
2a_{\mathrm{ter}}(\xi_\circ)(t_\circ+1)\ne0.
\]
Indeed, vanishing would give \(t_\circ=-1\); the constituent equality
would then give
\[
a_{\mathrm{bin}}(\xi_\circ)(1-p_\circ)^2=0,
\]
contrary to \(p_\circ<1\).
If the two constituent derivatives were equal, stationarity would make
this common derivative zero. Their difference is therefore nonzero, and
stationarity uniquely determines the mixture coordinate:
\[
u_\circ=(r_{\varepsilon_\circ}+1)\alpha_{\circ,S_6}
=
-\frac{\partial_t\Phi_{\mathrm{ter}}(\xi_\circ,t_\circ)}
       {\partial_t\Phi_{\mathrm{bin}}(\xi_\circ,t_\circ)
        -\partial_t\Phi_{\mathrm{ter}}(\xi_\circ,t_\circ)}.
\]
On a neighborhood where the derivative difference remains nonzero, set
\[
u_\sigma(\xi)
=
-\frac{\partial_t\Phi_{\mathrm{ter}}(\xi,t_\sigma(\xi))}
       {\partial_t\Phi_{\mathrm{bin}}(\xi,t_\sigma(\xi))
        -\partial_t\Phi_{\mathrm{ter}}(\xi,t_\sigma(\xi))}.
\]
It is continuous at \(\xi_\circ\), equals \(u_\circ\) there, and remains in
\((0,1)\) after shrinking the neighborhood. These arguments also apply
when \(p_\circ=1/2\).

We next continue the dual certificate while preserving its value at the
original point. Define
\[
\mathsf D_\varepsilon=
\begin{pmatrix}
r_\varepsilon&r_\varepsilon&1&1\\
1&1&r_\varepsilon&1\\
1&r_\varepsilon&r_\varepsilon&1\\
1&1&1&r_\varepsilon
\end{pmatrix},
\qquad
\boldsymbol f_{\mathrm{act}}(\xi)=
\begin{pmatrix}
f_{S_3}(\xi,t_\sigma(\xi))\\
f_{S_4}(\xi,t_\sigma(\xi))\\
f_{S_6}(\xi,t_\sigma(\xi))\\
f_{S_8}(\xi,t_\sigma(\xi))
\end{pmatrix}.
\]
For \(\varepsilon>0\), an explicit right inverse is
\[
\mathsf N_\varepsilon=
\begin{pmatrix}
\dfrac{r_\varepsilon+1}{d_\varepsilon(r_\varepsilon+2)}
&\dfrac{r_\varepsilon+1}{d_\varepsilon(r_\varepsilon+2)}
&-\dfrac1{d_\varepsilon}
&-\dfrac1{d_\varepsilon(r_\varepsilon+2)}\\
0&-\dfrac1{d_\varepsilon}&\dfrac1{d_\varepsilon}&0\\
-\dfrac1{d_\varepsilon(r_\varepsilon+2)}
&\dfrac{r_\varepsilon+1}{d_\varepsilon(r_\varepsilon+2)}
&0&-\dfrac1{d_\varepsilon(r_\varepsilon+2)}\\
-\dfrac1{d_\varepsilon(r_\varepsilon+2)}
&-\dfrac1{d_\varepsilon(r_\varepsilon+2)}
&0&\dfrac{r_\varepsilon+1}{d_\varepsilon(r_\varepsilon+2)}
\end{pmatrix}.
\]
Direct multiplication gives
\[
\mathsf D_\varepsilon\mathsf N_\varepsilon=\operatorname{Id}_4.
\]
Set
\[
\mathsf N_\circ=\mathsf N_{\varepsilon_\circ}.
\]
At \(\xi_\circ\), the matrix
\(\mathsf D_{\varepsilon_\circ}\mathsf N_\circ\) is the identity.
Its determinant remains nonzero nearby. On that neighborhood define
\[
\eta(\xi)
=
\eta_\circ+
\mathsf N_\circ
(\mathsf D_\varepsilon\mathsf N_\circ)^{-1}
\bigl(\boldsymbol f_{\mathrm{act}}(\xi)
      -\mathsf D_\varepsilon\eta_\circ\bigr).
\]
Matrix inversion is continuous there. The original active equalities
make the parenthesized residual zero at \(\xi_\circ\), and multiplication
by \(\mathsf D_\varepsilon\) gives
\[
\eta(\xi_\circ)=\eta_\circ,
\qquad
\mathsf D_\varepsilon\eta(\xi)=\boldsymbol f_{\mathrm{act}}(\xi).
\]
Thus the four active equalities indexed by \(S_3,S_4,S_6,S_8\) continue.

The fifth equality follows from the constituent equality. Indeed,
\[
\frac{b_{S_6}^\top\eta+b_{S_9}^\top\eta}{r_\varepsilon+1}
=
\frac{b_{S_3}^\top\eta+b_{S_4}^\top\eta+b_{S_8}^\top\eta}
     {r_\varepsilon+2},
\]
because both sides equal \(\mathbf1^\top\eta\).
Subtracting this identity from
\(\Phi_{\mathrm{bin}}(\xi,t_\sigma)=
\Phi_{\mathrm{ter}}(\xi,t_\sigma)\), using the four continued
equalities, gives
\[
b_{S_9}(\varepsilon)^\top\eta(\xi)
=
f_{S_9}(\xi,t_\sigma(\xi)).
\]

All pattern masses stay positive nearby. Hence the finitely many
information functions, projected scores, and dual-ray values are
continuous at \(\xi_\circ\). The strict inactive slacks and nonzero score
differences therefore persist simultaneously:
\[
\begin{gathered}
f_S(\xi,t_\sigma(\xi))
<
b_S(\varepsilon)^\top\eta(\xi)
\quad(S\notin\mathcal K_5),\\
\rho_S(\xi,t_\sigma(\xi))
\ne\rho_T(\xi,t_\sigma(\xi))
\quad(S,T\in\mathcal K_5,\ S\ne T).
\end{gathered}
\]
The strict interior inequalities, the branch denominator, and the
derivative difference likewise persist. For every point in the resulting
neighborhood, choose
\[
\alpha(\xi)=\alpha^{u_\sigma(\xi)}(\varepsilon).
\]
The normalization calculation gives feasibility and five positive
weights, while the defining formula for \(u_\sigma\) gives
\[
\partial_tF_{\theta(\xi)}
\bigl(\alpha(\xi),t_\sigma(\xi)\bigr)=0.
\]
Together with the continued dual equalities, strict slacks, and score
separation, these are all the certificate conditions. Every point of
\(\mathcal R_5\) consequently has a neighborhood in \(\mathbb R^4\)
contained in \(\mathcal R_5\), proving openness.
% lean: CausalSmith.Stat.LdpAteEfficiencySurface.isOpen_R5_continuousCertificate

\item \textit{The centered certificate and nonemptiness.}
Fix \(0<p<1/2\), and define
\[
\theta_{\mathrm{cen}}=\begin{pmatrix}1/2\\1/2\end{pmatrix},
\qquad
\xi_{\mathrm{cen}}(p)=(p,1/2,1/2,\log3),
\]
\[
\begin{aligned}
s_{\mathrm{cen}}&=\sqrt{\frac{8p^2}{5(1+p)}},
&
D_{\mathrm{cen}}&=1-2p+s_{\mathrm{cen}},\\
t_{\mathrm{cen}}&=\frac{p-s_{\mathrm{cen}}}{D_{\mathrm{cen}}},
&
\lambda_{\mathrm{cen}}&=\frac{s_{\mathrm{cen}}}{D_{\mathrm{cen}}},\\
J_{\mathrm{cen}}&=\frac{s_{\mathrm{cen}}^2(1-p)^2}{D_{\mathrm{cen}}^2},
&
\eta_{\mathrm{cen}}&=\frac{J_{\mathrm{cen}}}{4}(-1,-1,3,3)^\top.
\end{aligned}
\]
The identity
\[
5(1+p)s_{\mathrm{cen}}^2=8p^2
\]
and \(0<p<1/2\) give
\[
p<s_{\mathrm{cen}}<2p,\qquad
D_{\mathrm{cen}}>0,\qquad
0<\lambda_{\mathrm{cen}}<1,\qquad
t_{\mathrm{cen}}<0.
\]
Set
\[
\alpha_{\mathrm{cen}}=\alpha^{\lambda_{\mathrm{cen}}}(\log3).
\]
Thus
\[
\alpha_{\mathrm{cen},S_6}
=\alpha_{\mathrm{cen},S_9}=\frac{\lambda_{\mathrm{cen}}}{4},
\qquad
\alpha_{\mathrm{cen},S_3}
=\alpha_{\mathrm{cen},S_4}
=\alpha_{\mathrm{cen},S_8}
=\frac{1-\lambda_{\mathrm{cen}}}{5},
\]
and all other weights vanish. Its row sums are
\[
5\,\frac{1-\lambda_{\mathrm{cen}}}{5}
+
4\,\frac{\lambda_{\mathrm{cen}}}{4}=1,
\]
so it is feasible and has precisely the five stated positive weights.

At these parameters the constituent objectives reduce to
\[
\Phi_{\mathrm{bin}}(\xi_{\mathrm{cen}}(p),t)
=\{p-(1-2p)t\}^2,
\qquad
\Phi_{\mathrm{ter}}(\xi_{\mathrm{cen}}(p),t)
=s_{\mathrm{cen}}^2(t+1)^2.
\]
Define the positive score magnitudes
\[
\rho_{\mathrm{bin}}
=p-(1-2p)t_{\mathrm{cen}}
=\frac{s_{\mathrm{cen}}q}{D_{\mathrm{cen}}},
\qquad
\rho_{\mathrm{ter}}
=\frac{2p(t_{\mathrm{cen}}+1)}{1+p}.
\]
Direct rearrangement gives
\[
t_{\mathrm{cen}}+1=\frac{q}{D_{\mathrm{cen}}},
\qquad
s_{\mathrm{cen}}(t_{\mathrm{cen}}+1)=\rho_{\mathrm{bin}},
\qquad
J_{\mathrm{cen}}=\rho_{\mathrm{bin}}^2>0.
\]
Consequently,
\[
\begin{aligned}
\partial_tF_{\theta_{\mathrm{cen}}}
(\alpha_{\mathrm{cen}},t_{\mathrm{cen}})
&=-2\lambda_{\mathrm{cen}}(1-2p)\rho_{\mathrm{bin}}
  +2(1-\lambda_{\mathrm{cen}})s_{\mathrm{cen}}\rho_{\mathrm{bin}}\\
&=0.
\end{aligned}
\]

For completeness, the information and dual calculations over all
fourteen patterns are
\[
\begin{array}{c|c|c}
\text{patterns}
&b_S(\log3)^\top\eta_{\mathrm{cen}}/J_{\mathrm{cen}}
&f_S(\xi_{\mathrm{cen}}(p),t_{\mathrm{cen}})/J_{\mathrm{cen}}\\ \hline
S_1,S_2&1/2&
\dfrac{4q^2t_{\mathrm{cen}}^2}{(2-p)J_{\mathrm{cen}}}\\
S_3&0&0\\
S_4,S_8&5/2&5/2\\
S_5,S_{10}&2&
\dfrac{2(p+t_{\mathrm{cen}})^2}{J_{\mathrm{cen}}}\\
S_6,S_9&2&2\\
S_7,S_{11}&3/2&
\dfrac{5(1+p)}{2(3-p)}\\
S_{12}&4&0\\
S_{13},S_{14}&7/2&
\dfrac{4q^2t_{\mathrm{cen}}^2}{(2+p)J_{\mathrm{cen}}}
\end{array}
\]
This table follows by inserting the displayed record probabilities and
binary-mask enumeration into \(h_S\), \(g_S\), and \(b_S^\top\eta_{\mathrm{cen}}\).
For example, for \(S_4=\{x_2\}\),
\[
h_{S_4}(\xi_{\mathrm{cen}}(p))=1+p,\qquad
g_{S_4}(\xi_{\mathrm{cen}}(p))=(0,-2p)^\top,
\]
so
\[
f_{S_4}(\xi_{\mathrm{cen}}(p),t_{\mathrm{cen}})
=
\frac{4p^2(t_{\mathrm{cen}}+1)^2}{1+p}
=\frac52J_{\mathrm{cen}},
\]
and its ray gives the same dual value.

We verify strictness in the five inactive families. First,
\[
5s_{\mathrm{cen}}^2<8p^2
\quad\Longrightarrow\quad
s_{\mathrm{cen}}<\frac{4p}{3}.
\]
Together with \(p<s_{\mathrm{cen}}\), this yields
\[
8(p-s_{\mathrm{cen}})^2
<
\frac89p^2
<
\frac32p^2
<
(2-p)s_{\mathrm{cen}}^2.
\]
Multiplying by \(q^2/D_{\mathrm{cen}}^2>0\) and dividing by \(2(2-p)\)
proves
\[
\frac{4q^2t_{\mathrm{cen}}^2}{2-p}<\frac12J_{\mathrm{cen}},
\]
the strict slack for \(S_1,S_2\).

Next,
\[
p+t_{\mathrm{cen}}
=\frac{q(2p-s_{\mathrm{cen}})}{D_{\mathrm{cen}}},
\qquad
0<p+t_{\mathrm{cen}}<\frac{qs_{\mathrm{cen}}}{D_{\mathrm{cen}}}
=\rho_{\mathrm{bin}}.
\]
Squaring proves
\[
2(p+t_{\mathrm{cen}})^2<2J_{\mathrm{cen}},
\]
which handles \(S_5,S_{10}\).

For \(S_7,S_{11}\), the same square identity gives
\[
3s_{\mathrm{cen}}^2(3-p)-8p^2
=
\frac{8p^2(4-8p)}{5(1+p)}>0.
\]
Using \(t_{\mathrm{cen}}+1=q/D_{\mathrm{cen}}\), this is exactly
\[
\frac{4p^2(t_{\mathrm{cen}}+1)^2}{3-p}
<
\frac32J_{\mathrm{cen}}.
\]
For \(S_{12}\), strictness is \(0<4J_{\mathrm{cen}}\).
Finally,
\[
8(p-s_{\mathrm{cen}})^2
<
(2-p)s_{\mathrm{cen}}^2
<
7(2+p)s_{\mathrm{cen}}^2
\]
gives
\[
\frac{4q^2t_{\mathrm{cen}}^2}{2+p}
<
\frac72J_{\mathrm{cen}},
\]
as required for \(S_{13},S_{14}\).

The active projected scores, in the order \(S_3,S_4,S_6,S_8,S_9\), are
\[
\bigl(0,-\rho_{\mathrm{ter}},-\rho_{\mathrm{bin}},
          \rho_{\mathrm{ter}},\rho_{\mathrm{bin}}\bigr).
\]
Both magnitudes are positive, and
\[
\left(\frac{\rho_{\mathrm{ter}}}{\rho_{\mathrm{bin}}}\right)^2
=
\frac{5}{2(1+p)}>1.
\]
Hence all five scores are distinct: the comparisons between positive
magnitudes are strict, and the remaining comparisons follow from their
signs. Reversing the order of a pair gives the reverse comparison.

We have verified feasibility, positive active weights, stationarity,
all active dual equalities, all strict inactive slacks, and score
separation. Therefore
\[
\xi_{\mathrm{cen}}(p)\in\mathcal R_5
\qquad(0<p<1/2).
\]
Taking \(p=1/4\) proves that \(\mathcal R_5\) is nonempty.
% lean: CausalSmith.Stat.LdpAteEfficiencySurface.R5_nonempty_centered

\item \textit{Centered optimality and uniqueness.}
Stationarity and the nonnegative quadratic coefficients give
\[
\begin{aligned}
&F_{\theta_{\mathrm{cen}}}(\alpha_{\mathrm{cen}},t)
-F_{\theta_{\mathrm{cen}}}(\alpha_{\mathrm{cen}},t_{\mathrm{cen}})\\
&\qquad=
\bigl\{\lambda_{\mathrm{cen}}(1-2p)^2
 +(1-\lambda_{\mathrm{cen}})s_{\mathrm{cen}}^2\bigr\}
(t-t_{\mathrm{cen}})^2\ge0.
\end{aligned}
\]
Its value at \(t_{\mathrm{cen}}\) is \(J_{\mathrm{cen}}\).
For every feasible \(\beta\), nonnegativity of its weights and the dual
bounds give
\[
\begin{aligned}
\inf_{t\in\mathbb R}F_{\theta_{\mathrm{cen}}}(\beta,t)
&\le F_{\theta_{\mathrm{cen}}}(\beta,t_{\mathrm{cen}})\\
&\le\sum_S\beta_S b_S(\log3)^\top\eta_{\mathrm{cen}}\\
&=(B_{\log3}\beta)^\top\eta_{\mathrm{cen}}
=\mathbf1^\top\eta_{\mathrm{cen}}
=J_{\mathrm{cen}}.
\end{aligned}
\]
The infima here are bounded below by zero and are taken over nonempty
sets. The feasible vector \(\alpha_{\mathrm{cen}}\) reaches the upper
bound, so the optimization in \cref{obj:def:staircase-saddle} yields
\[
J^*(\theta_{\mathrm{cen}},p,\log3)
=
J_{\mathrm{cen}}
=
\inf_tF_{\theta_{\mathrm{cen}}}(\alpha_{\mathrm{cen}},t).
\]

Suppose a feasible \(\beta\) has this optimized profile. Every inequality
in the preceding bound is then an equality. In particular,
\[
\sum_S\beta_S
\bigl\{b_S(\log3)^\top\eta_{\mathrm{cen}}
       -f_S(\xi_{\mathrm{cen}}(p),t_{\mathrm{cen}})\bigr\}=0.
\]
Each summand is nonnegative, and every inactive slack is strictly
positive. Thus
\[
\beta_S=0\qquad(S\notin\mathcal K_5).
\]
Also,
\[
F_{\theta_{\mathrm{cen}}}(\beta,t_{\mathrm{cen}})
=
J_{\mathrm{cen}}
\le F_{\theta_{\mathrm{cen}}}(\beta,t)
\quad(t\in\mathbb R),
\]
so differentiation of this quadratic at its minimum gives
\[
\partial_tF_{\theta_{\mathrm{cen}}}(\beta,t_{\mathrm{cen}})=0.
\]

The row constraints already show that
\[
\beta_{S_3}=\beta_{S_4}=\beta_{S_8}
=\frac{1-4\beta_{S_6}}5,
\qquad
\beta_{S_9}=\beta_{S_6}.
\]
To solve the remaining stationarity equation, define its reduced
coefficient
\[
C_{\mathrm{cen}}
=
10s_{\mathrm{cen}}p^2+5s_{\mathrm{cen}}p
-5s_{\mathrm{cen}}-8p^2.
\]
Substitute these row coordinates into the displayed derivative formula,
use \(t_{\mathrm{cen}}=(p-s_{\mathrm{cen}})/D_{\mathrm{cen}}\), and clear
the positive denominators. The resulting equation is
\[
(p-1)\bigl(C_{\mathrm{cen}}\beta_{S_6}+2p^2\bigr)=0.
\]
Since \(p<1/2\), it follows that
\[
C_{\mathrm{cen}}\beta_{S_6}=-2p^2.
\]
The square identity for \(s_{\mathrm{cen}}\) further gives
\[
\begin{aligned}
C_{\mathrm{cen}}s_{\mathrm{cen}}
&=5(2p-1)(1+p)s_{\mathrm{cen}}^2-8p^2s_{\mathrm{cen}}\\
&=-8p^2D_{\mathrm{cen}}<0.
\end{aligned}
\]
Thus \(C_{\mathrm{cen}}\ne0\), and
\[
C_{\mathrm{cen}}\bigl(4D_{\mathrm{cen}}\beta_{S_6}\bigr)
=
-8p^2D_{\mathrm{cen}}
=
C_{\mathrm{cen}}s_{\mathrm{cen}}.
\]
Cancellation proves
\[
\beta_{S_6}=\frac{s_{\mathrm{cen}}}{4D_{\mathrm{cen}}}
=\alpha_{\mathrm{cen},S_6}.
\]
The row-coordinate identities determine the other four active weights,
and both vectors vanish elsewhere. Hence
\[
\beta=\alpha_{\mathrm{cen}},
\]
proving the centered uniqueness assertion.
% lean: CausalSmith.Stat.LdpAteEfficiencySurface.centeredR5_profile_optimal_unique

\item \textit{Attainment, refinement, and the centered cardinality bound.}
We first give the directional inequality used for attaining channels.
For every \(\xi\in\mathcal U\), define the singleton design
\[
\alpha^{\mathrm{sgl}}_S=
\begin{cases}
1/(r_\varepsilon+3),&S\in\{S_1,S_2,S_4,S_8\},\\
0,&\text{otherwise}.
\end{cases}
\]
These four patterns are the four singleton subsets. Their rays sum to
\((r_\varepsilon+3)\mathbf1\), proving feasibility. Its objective is a
quadratic with nonnegative terms, and in particular
\[
F_\theta(\alpha^{\mathrm{sgl}},t)
\ge
\frac{d_\varepsilon^2q^2}
     {(r_\varepsilon+3)h_{S_1}(\xi)}t^2,
\qquad
F_\theta(\alpha^{\mathrm{sgl}},t)
\ge
\frac{d_\varepsilon^2p^2}
     {(r_\varepsilon+3)h_{S_8}(\xi)}(t+1)^2.
\]
Both coefficients are positive. This quadratic has a minimizer, and a
zero minimum would force both \(t=0\) and \(t+1=0\). Therefore
\[
0<\inf_tF_\theta(\alpha^{\mathrm{sgl}},t)
\le J^*(\theta,p,\varepsilon),
\]
where the final comparison is supplied by the optimization in
\cref{obj:def:staircase-saddle}.

For a stationary private channel \(Q\), write
\[
\mathsf P_Q=P_\theta Q,\qquad
\nu_Q=\sum_{j=0}^3Q(x_j,\cdot),
\]
and define its output score by
\[
s_Q(z)=
\begin{cases}
\displaystyle
\frac{\nabla_\theta\{d\mathsf P_Q/d\nu_Q\}(z)}
     {\{d\mathsf P_Q/d\nu_Q\}(z)},
&\{d\mathsf P_Q/d\nu_Q\}(z)>0,\\[6pt]
0,&\{d\mathsf P_Q/d\nu_Q\}(z)=0.
\end{cases}
\]
The score representation of the information matrix in
\cref{obj:lem:staircase-refinement} gives, with
\(I_Q=I_\theta(Q)\),
\[
I_Q=\int s_Q(z)s_Q(z)^\top\,d\mathsf P_Q(z),
\qquad
v_{\mathrm{sol}}^\top I_Qv_{\mathrm{sol}}
=
\int(v_{\mathrm{sol}}^\top s_Q(z))^2\,d\mathsf P_Q(z)\ge0
\quad(v_{\mathrm{sol}}\in\mathbb R^2).
\]
Thus \(I_Q\) is positive semidefinite.

Let \(\mathcal V(I)\) denote the extended contrast variance:
\[
\mathcal V(I)=
\begin{cases}
\displaystyle
\max\!\left\{\inf_{\{v_{\mathrm{sol}}\in\mathbb R^2:
                         Iv_{\mathrm{sol}}=c\}}
                  c^\top v_{\mathrm{sol}},\,0\right\},
&c\in\operatorname{range}(I),\\
+\infty,&c\notin\operatorname{range}(I),
\end{cases}
\qquad
c=\begin{pmatrix}-1\\1\end{pmatrix}.
\]
If \(Q\) attains the variance in the statement, its variance is finite,
so there is a vector \(v_{\mathrm{sol}}\in\mathbb R^2\) satisfying
\[
I_Qv_{\mathrm{sol}}=c.
\]
Define
\[
V_{\mathrm{sol}}=c^\top v_{\mathrm{sol}}
=v_{\mathrm{sol}}^\top I_Qv_{\mathrm{sol}}>0.
\]
To justify strict positivity, zero quadratic value for a positive
semidefinite matrix implies \(I_Qv_{\mathrm{sol}}=0\), contradicting
\(c\ne0\). Any other solution has the same contrast pairing, since
\[
I_Qv_{\mathrm{sol}}'=c
\quad\Longrightarrow\quad
c^\top v_{\mathrm{sol}}'
=
v_{\mathrm{sol}}^\top I_Qv_{\mathrm{sol}}'
=
v_{\mathrm{sol}}^\top c
=
V_{\mathrm{sol}}.
\]
Consequently attainment and positivity of \(J^*\) imply
\[
V_{\mathrm{sol}}=\mathcal V(I_Q)
=V^*(\theta,p,\varepsilon)
=\frac1{J^*(\theta,p,\varepsilon)}.
\]
Here both values are positive, so the nonnegative-part convention
preserves their real values.

For arbitrary \(t\in\mathbb R\), define the residual vector
\[
v_{\mathrm{res}}=V_{\mathrm{sol}}v(t)-v_{\mathrm{sol}}.
\]
Since \(c^\top v(t)=1\), positive semidefiniteness gives
\[
\begin{aligned}
0
&\le v_{\mathrm{res}}^\top I_Qv_{\mathrm{res}}\\
&=V_{\mathrm{sol}}^2v(t)^\top I_Qv(t)-V_{\mathrm{sol}}.
\end{aligned}
\]
Dividing by \(V_{\mathrm{sol}}>0\) and substituting its value proves
\[
J^*(\theta,p,\varepsilon)\le v(t)^\top I_Qv(t)
\qquad(t\in\mathbb R).
\]
% lean: CausalSmith.Stat.LdpAteEfficiencySurface.attainment_forces_direction_lower

Return to the centered parameters and an attaining \(Q\). Use the
constant privacy sequence \(\varepsilon_n=\log3\); it satisfies
\cref{obj:ass:fixed-privacy}. The centered parameters also satisfy
\cref{obj:ass:interior-assignment,obj:ass:interior-means}.
By \cref{obj:lem:staircase-refinement}, there are feasible weights
\(\beta\) and a probability kernel \(K\) such that
\[
Q=K\circ Q_\beta,
\qquad
I_{\theta_{\mathrm{cen}}}(\beta)-I_{\theta_{\mathrm{cen}}}(Q)\succeq0.
\]
For every \(t\), the directional lower bound and this matrix order imply
\[
J_{\mathrm{cen}}
\le
v(t)^\top I_{\theta_{\mathrm{cen}}}(Q)v(t)
\le
v(t)^\top I_{\theta_{\mathrm{cen}}}(\beta)v(t)
=
F_{\theta_{\mathrm{cen}}}(\beta,t).
\]
The centered dual bound supplies the reverse inequality at
\(t_{\mathrm{cen}}\). Hence
\[
F_{\theta_{\mathrm{cen}}}(\beta,t_{\mathrm{cen}})
=J_{\mathrm{cen}}.
\]
Equality in the strict dual bound forces \(\beta\) to vanish outside
\(\mathcal K_5\). The preceding inequality for every \(t\) makes
\(t_{\mathrm{cen}}\) a minimizer for \(\beta\), so its derivative there is
zero. The supported stationary calculation above therefore gives
\[
\beta=\alpha_{\mathrm{cen}},
\qquad
\beta_S>0\quad(S\in\mathcal K_5).
\]
Moreover, the directional inequalities at \(t_{\mathrm{cen}}\) give
\[
v(t_{\mathrm{cen}})^\top I_{\theta_{\mathrm{cen}}}(\beta)
v(t_{\mathrm{cen}})
=
v(t_{\mathrm{cen}})^\top I_{\theta_{\mathrm{cen}}}(Q)
v(t_{\mathrm{cen}})
=
J_{\mathrm{cen}}.
\]
The score-equality conclusion of
\cref{obj:lem:staircase-refinement} now applies. If two distinct active
conditional output measures were not mutually singular, their projected
scores would agree. The centered score calculation proves that they
differ. Thus, writing \(K_S=K(S,\cdot)\),
\[
K_S\perp K_T
\qquad(S,T\in\mathcal K_5,\ S\ne T).
\]

We spell out how these five singular probability measures give the
cardinality bound. Evaluations of output-space measures on arbitrary
subsets use their outer-measure extensions. For any probability measure
\(\mathsf P_{\mathrm{out}}\) on \(Z\), define
\[
\mathsf P_{\mathrm{out}}^{\mathrm{out}}(A)
=
\inf\bigl\{\mathsf P_{\mathrm{out}}(E):
A\subseteq E\subseteq Z,\ E\text{ measurable}\bigr\},
\qquad A\subseteq Z.
\]
We also write \(\mathsf P_{\mathrm{out}}(A)\) for this extension; it agrees
with the original measure on measurable sets. In particular, singleton
evaluations of \(K_S\) and \(\mathsf P_Q\) below use this convention, as
does the output cardinality:
\[
\kappa(Q)=
\begin{cases}
\#\{z\in Z:\mathsf P_Q^{\mathrm{out}}(\{z\})\ne0\},&Z\text{ finite},\\
+\infty,&Z\text{ infinite}.
\end{cases}
\]
These extensions are monotone and finitely subadditive. Absolute
continuity transfers to their evaluations on arbitrary sets: a set of
outer measure zero is contained in a measurable null set, obtained by
intersecting measurable supersets whose masses tend to zero; the
dominated measure also vanishes on that measurable null set and hence
on the original set.

For every measurable output event \(A\), factorization
and the positive record probability \(\pi_{\theta,0}>0\) give
\[
\beta_S(b_S(\varepsilon))_0K_S(A)\le Q(x_0,A),
\qquad
\pi_{\theta,0}Q(x_0,A)\le\mathsf P_Q(A).
\]
The coefficients on the left are positive for active \(S\). Therefore
\[
K_S\ll Q(x_0,\cdot)\ll\mathsf P_Q
\qquad(S\in\mathcal K_5).
\]
Suppose first that the output space \(Z\) is finite. For each
\(S\in\mathcal K_5\), choose a point \(z_S\in Z\) such that
\[
K_S(\{z_S\})>0.
\]
Such a point exists because otherwise finite subadditivity would give
\[
1=K_S(Z)
=K_S\!\left(\bigcup_{z\in Z}\{z\}\right)
\le\sum_{z\in Z}K_S(\{z\})=0.
\]
This contradiction also covers an empty output space, with the sum
interpreted as zero. Absolute continuity then gives
\[
\mathsf P_Q(\{z_S\})>0\qquad(S\in\mathcal K_5),
\]
since a zero output-law mass would force \(K_S(\{z_S\})=0\).

To prove that the chosen points are distinct, fix distinct
\(S,T\in\mathcal K_5\). Mutual singularity supplies a measurable set
\(E_{S,T}\subseteq Z\) with
\[
K_S(E_{S,T})=0,\qquad K_T(Z\setminus E_{S,T})=0.
\]
Suppose \(z_S=z_T\). If \(z_S\in E_{S,T}\), measure monotonicity gives
\[
0<K_S(\{z_S\})\le K_S(E_{S,T})=0.
\]
If \(z_S\notin E_{S,T}\), then \(z_T\in Z\setminus E_{S,T}\), and
\[
0<K_T(\{z_T\})\le K_T(Z\setminus E_{S,T})=0.
\]
Both cases are contradictory. Thus the five selected points are distinct
and have positive output-law mass, proving
\[
5=|\mathcal K_5|
=\#\{z_S:S\in\mathcal K_5\}
\le
\#\{z\in Z:\mathsf P_Q(\{z\})\ne0\}
=\kappa(Q).
\]
For an infinite output space, the stated convention gives
\(\kappa(Q)=+\infty\), and the same bound follows.
% lean: CausalSmith.Stat.LdpAteEfficiencySurface.centeredR5_outputCardinality_ge_five

\item \textit{Exchange of the treatment labels.}
Suppose \(1/2<p<1\) and put \(q=1-p\). Then
\[
0<q<1/2,
\qquad
(q,1/2,1/2,\log3)\in\mathcal R_5.
\]
Define the input exchange, its action on patterns and weights, and the
coordinate-exchange matrix by
\[
\begin{gathered}
\mathfrak s(x_0)=x_2,\quad
\mathfrak s(x_1)=x_3,\quad
\mathfrak s(x_2)=x_0,\quad
\mathfrak s(x_3)=x_1,\\
\mathfrak s S=\{\mathfrak s x:x\in S\},\qquad
\alpha^{\mathrm{sw}}_S=\alpha_{\mathfrak s S},\qquad
\Pi=\begin{pmatrix}0&1\\1&0\end{pmatrix}.
\end{gathered}
\]
Both exchanges are involutions. The ray identity
\[
(b_{\mathfrak s S}(\varepsilon))_j
=b_S(\varepsilon)_{\mathfrak s(j)}
\]
shows, by reindexing the normalization sums, that
\[
\alpha\in\mathcal A_\varepsilon
\quad\Longleftrightarrow\quad
\alpha^{\mathrm{sw}}\in\mathcal A_\varepsilon.
\]
On the centered slice, exchanging the record probabilities gives
\[
h_{\mathfrak s S}(\xi_{\mathrm{cen}}(p))
=h_S(\xi_{\mathrm{cen}}(q)),
\qquad
g_{\mathfrak s S}(\xi_{\mathrm{cen}}(p))
=\Pi g_S(\xi_{\mathrm{cen}}(q)).
\]
Since
\[
\Pi v(t)=-v(-1-t),
\]
the information contributions satisfy
\[
f_{\mathfrak s S}(\xi_{\mathrm{cen}}(p),t)
=
f_S(\xi_{\mathrm{cen}}(q),-1-t).
\]
Reindexing the sum and using the bijection \(t\mapsto-1-t\) consequently
gives
\[
\inf_t\sum_S\alpha^{\mathrm{sw}}_S
 f_S(\xi_{\mathrm{cen}}(p),t)
=
\inf_t\sum_S\alpha_S f_S(\xi_{\mathrm{cen}}(q),t).
\]
The exchange is also a bijection of feasible weights. Taking the
optimized profiles on both sides proves
\[
J^*(\theta_{\mathrm{cen}},p,\log3)
=
J^*(\theta_{\mathrm{cen}},q,\log3),
\qquad
V^*(\theta_{\mathrm{cen}},p,\log3)
=
V^*(\theta_{\mathrm{cen}},q,\log3).
\]

Choose the centered optimizer at \(q\) and exchange its weights. This
produces a feasible optimizer at \(p\) with support
\[
\mathfrak s\mathcal K_5=\{\mathfrak s S:S\in\mathcal K_5\}.
\]
If another feasible vector \(\beta\) is optimal at \(p\), then
\(\beta^{\mathrm{sw}}\) is optimal at \(q\). Centered uniqueness gives
\[
\beta^{\mathrm{sw}}=\alpha_{\mathrm{cen}}(q).
\]
Applying the involution again proves
\[
\beta=\alpha_{\mathrm{cen}}(q)^{\mathrm{sw}},
\]
the claimed uniqueness at the original assignment probability.

For the channel assertion, define
\[
Q^{\mathrm{sw}}(x,A)=Q(\mathfrak s x,A).
\]
This is a probability kernel and satisfies the same privacy inequalities
as \(Q\), since those inequalities hold for every input pair.
In the next identities, quantities for \(Q^{\mathrm{sw}}\) are evaluated
at assignment probability \(q\), and those for \(Q\) at \(p\).
Reindexing the four record probabilities gives equality of the output
laws; reindexing their derivatives exchanges the two parameter
coordinates. Therefore
\[
P_{\theta_{\mathrm{cen}}}Q^{\mathrm{sw}}
=
P_{\theta_{\mathrm{cen}}}Q,
\qquad
I_{\theta_{\mathrm{cen}}}(Q^{\mathrm{sw}})
=
\Pi I_{\theta_{\mathrm{cen}}}(Q)\Pi.
\]
The contrast variance is preserved by this matrix exchange. Indeed,
\(\Pi c=-c\), and for every \(v_{\mathrm{sol}}\in\mathbb R^2\),
\[
Iv_{\mathrm{sol}}=c
\quad\Longleftrightarrow\quad
(\Pi I\Pi)(-\Pi v_{\mathrm{sol}})=c,
\qquad
c^\top(-\Pi v_{\mathrm{sol}})=c^\top v_{\mathrm{sol}}.
\]
Thus the solution-value sets coincide; if they are empty, both contrast
variances equal \(+\infty\). In either case,
\[
\mathcal V(\Pi I\Pi)=\mathcal V(I).
\]
An attaining \(Q\) at \(p\) consequently gives an attaining
\(Q^{\mathrm{sw}}\) at \(q\). The centered bound yields
\(\kappa(Q^{\mathrm{sw}})\ge5\). Equality of their output laws, on the same
output space, gives
\[
\kappa(Q^{\mathrm{sw}})=\kappa(Q),
\]
and hence the desired bound at \(p\).
% lean: CausalSmith.Stat.LdpAteEfficiencySurface.reflected_five_output_region

\item \textit{Uniqueness for a general five-pattern certificate.}
Fix \(\xi=(p,\mu_0,\mu_1,\varepsilon)\in\mathcal R_5\), put
\(\theta=\theta(\xi)\), and choose a certificate \((\alpha_0,t_0,\eta)\).
For any feasible weight vector, its objective is nonnegative.
For the certificate vector, expansion of the quadratic gives
\[
\begin{aligned}
F_\theta(\alpha_0,t)-F_\theta(\alpha_0,t_0)
&=(t-t_0)\partial_tF_\theta(\alpha_0,t_0)\\
&\quad+(t-t_0)^2
 \sum_S\alpha_{0,S}
 \frac{(g_S(\xi)^\top e_{\mathrm{prof}})^2}{h_S(\xi)}\\
&\ge0,
\end{aligned}
\]
where stationarity makes the linear term zero.
The active equalities, the vanishing inactive weights, and normalization
give
\[
F_\theta(\alpha_0,t_0)
=
\sum_S\alpha_{0,S}b_S(\varepsilon)^\top\eta
=
\sum_{j=0}^3\eta_j.
\]
For every feasible \(\beta\),
\[
\inf_tF_\theta(\beta,t)
\le F_\theta(\beta,t_0)
\le\sum_S\beta_Sb_S(\varepsilon)^\top\eta
=\sum_{j=0}^3\eta_j.
\]
The certificate vector attains this upper bound. Hence
\[
J^*(\theta,p,\varepsilon)
=
\inf_tF_\theta(\alpha_0,t)
=
F_\theta(\alpha_0,t_0)
=
\sum_{j=0}^3\eta_j.
\]
Moreover, equality of another feasible profile with \(J^*\) forces
equality in the dual bound and therefore
\[
\beta_S=0\qquad(S\notin\mathcal K_5),
\]
by the same nonnegative slack-sum argument.

The constituent coefficients are positive, and the active dual
equalities give
\[
a_{\mathrm{bin}}(\xi)\{p-(1-2p)t_0\}^2
=
a_{\mathrm{ter}}(\xi)(t_0+1)^2.
\]
As already derived, \(t_0=-1\) would force \(p=1\), so
\[
\partial_t\Phi_{\mathrm{ter}}(\xi,t_0)\ne0.
\]
Define the certificate mixture coordinate by
\[
u_0=(r_\varepsilon+1)\alpha_{0,S_6}.
\]
The normalization and stationarity calculations give
\[
u_0
=
-\frac{\partial_t\Phi_{\mathrm{ter}}(\xi,t_0)}
       {\partial_t\Phi_{\mathrm{bin}}(\xi,t_0)
        -\partial_t\Phi_{\mathrm{ter}}(\xi,t_0)},
\qquad
\partial_t\Phi_{\mathrm{bin}}(\xi,t_0)
-\partial_t\Phi_{\mathrm{ter}}(\xi,t_0)\ne0.
\]

Now let \(\beta\) be any feasible optimized vector. Its nonnegative
objective has profile \(J^*\), and the dual upper bound at \(t_0\) gives
\[
J^*\le F_\theta(\beta,t_0)\le J^*.
\]
It follows that
\[
F_\theta(\beta,t_0)=J^*
\le F_\theta(\beta,t)\quad(t\in\mathbb R),
\qquad
\partial_tF_\theta(\beta,t_0)=0.
\]
Since \(\beta\) vanishes outside \(\mathcal K_5\), define
\[
u_\beta=(r_\varepsilon+1)\beta_{S_6}.
\]
Its stationarity equation yields
\[
u_\beta
=
-\frac{\partial_t\Phi_{\mathrm{ter}}(\xi,t_0)}
       {\partial_t\Phi_{\mathrm{bin}}(\xi,t_0)
        -\partial_t\Phi_{\mathrm{ter}}(\xi,t_0)}
=u_0.
\]
The supported normalization formula determines every weight from this
coordinate. Hence \(\beta=\alpha_0\). In particular, the unique optimized
vector has positive weight on all five active patterns.
% lean: CausalSmith.Stat.LdpAteEfficiencySurface.r5Certificate_unique_optimum

\item \textit{The cardinality bound throughout the region.}
Keep the general certificate from the preceding argument and let \(Q\)
be any stationary \(\varepsilon\)-LDP channel attaining the specified
extended contrast variance. The certificate includes the interior
conditions and \(\varepsilon>0\). With the constant sequence
\(\varepsilon_n=\varepsilon\), all hypotheses of
\cref{obj:lem:staircase-refinement} hold. It supplies feasible weights
\(\beta\) and a probability kernel \(K\) with
\[
Q=K\circ Q_\beta,
\qquad
I_\theta(\beta)-I_\theta(Q)\succeq0.
\]
The directional attainment bound already derived applies to these
parameters. Consequently, for every \(t\in\mathbb R\),
\[
J^*(\theta,p,\varepsilon)
\le v(t)^\top I_\theta(Q)v(t)
\le v(t)^\top I_\theta(\beta)v(t)
=F_\theta(\beta,t).
\]
At the certificate direction, the dual bound is
\[
F_\theta(\beta,t_0)\le
\sum_{j=0}^3\eta_j
=
J^*(\theta,p,\varepsilon).
\]
Thus
\[
F_\theta(\beta,t_0)=J^*(\theta,p,\varepsilon)
\le F_\theta(\beta,t)\quad(t\in\mathbb R),
\]
and
\[
\inf_tF_\theta(\beta,t)=J^*(\theta,p,\varepsilon).
\]
The general certificate uniqueness just proved gives
\[
\beta=\alpha_0,\qquad \beta_S>0\quad(S\in\mathcal K_5).
\]
The same upper and lower bounds at \(t_0\) also give directional equality:
\[
v(t_0)^\top I_\theta(\beta)v(t_0)
=
v(t_0)^\top I_\theta(Q)v(t_0)
=
J^*(\theta,p,\varepsilon).
\]
By the score-equality conclusion of
\cref{obj:lem:staircase-refinement}, overlapping active conditional
measures would have equal projected scores. The region's certificate
separates every pair, so
\[
K(S,\cdot)\perp K(T,\cdot)
\qquad(S,T\in\mathcal K_5,\ S\ne T).
\]
Factorization, positive active weights, positive ray coordinates, and
\(\pi_{\theta,0}=q(1-\mu_0)>0\) give
\[
K(S,\cdot)\ll Q(x_0,\cdot)\ll P_\theta Q
\qquad(S\in\mathcal K_5).
\]
Suppose first that the output space \(Z\) is finite. With the outer-measure
convention above, the conditional-atom argument chooses, for every
\(S\in\mathcal K_5\), a point
\(z_S\in Z\) satisfying
\[
K(S,\{z_S\})>0.
\]
Indeed, if all singleton masses of this probability measure were zero,
finite subadditivity would give
\[
1=K(S,Z)\le\sum_{z\in Z}K(S,\{z\})=0,
\]
including the empty-space case. Absolute continuity transfers positivity:
\[
(P_\theta Q)(\{z_S\})>0\qquad(S\in\mathcal K_5).
\]
For distinct \(S,T\in\mathcal K_5\), choose a measurable separating set
\(E_{S,T}\subseteq Z\) such that
\[
K(S,E_{S,T})=0,\qquad K(T,Z\setminus E_{S,T})=0.
\]
If \(z_S=z_T\), membership in this set gives the two cases
\[
\begin{aligned}
z_S\in E_{S,T}
&\quad\Longrightarrow\quad
0<K(S,\{z_S\})\le K(S,E_{S,T})=0,\\
z_S\notin E_{S,T}
&\quad\Longrightarrow\quad
0<K(T,\{z_T\})\le K(T,Z\setminus E_{S,T})=0.
\end{aligned}
\]
Hence the selected points are distinct, and
\[
5=|\mathcal K_5|
=\#\{z_S:S\in\mathcal K_5\}
\le\#\{z\in Z:(P_\theta Q)(\{z\})\ne0\}
=\kappa(Q).
\]
For an infinite output space the cardinality convention gives
\(\kappa(Q)=+\infty\). Thus in both cases
\[
\kappa(Q)\ge|\mathcal K_5|=5.
\]
This proves the universal assertion throughout \(\mathcal R_5\), and
completes all the claims.
% lean: CausalSmith.Stat.LdpAteEfficiencySurface.r5Certificate_outputCardinality_ge_five
\end{enumerate}
% lean: CausalSmith.Stat.LdpAteEfficiencySurface.five_output_region
\end{proof}
